\subsection{屈内特（Künneth）公式}

在本节中，我们考虑一个右 A-模的复形 (C, d) 和一个左 A-模的复形 (C', d')。考虑典范正合序列

\[
0 \rightarrow \mathrm{Z} (\mathrm{C}) \stackrel {{j}} {{\rightarrow}} \mathrm{C} \stackrel {{\delta}} {{\rightarrow}} \mathrm{B} (\mathrm{C}) (- 1) \rightarrow 0,\tag{I}
\]

\[
0 \rightarrow \mathrm{B} (\mathrm{C}) \xrightarrow {i} \mathrm{Z} (\mathrm{C}) \xrightarrow {p} \mathrm{H} (\mathrm{C}) \rightarrow 0;\tag{II}
\]

我们从 δ 推出一个 k-同态

\[
\mathrm{H} (\delta \otimes 1): \mathrm{H} (\mathrm{C} \otimes_ {\mathrm{A}} \mathrm{C} ^ {\prime}) \rightarrow \mathrm{H} (\mathrm{B} (\mathrm{C}) \otimes_ {\mathrm{A}} \mathrm{C} ^ {\prime}) (- 1);
\]

我们由(II)推出一个连接同态

\[
\partial (\mathrm{II}, \mathrm{H} (\mathrm{C} ^ {\prime})): \operatorname{Tor} _ {1} ^ {\mathbf {A}} (\mathrm{H} (\mathrm{C}), \mathrm{H} (\mathrm{C} ^ {\prime})) \rightarrow \mathrm{B} (\mathrm{C}) \otimes_ {\mathbf {A}} \mathrm{H} (\mathrm{C} ^ {\prime});
\]

如果我们装备 \(\operatorname{Tor}_{1}^{\mathbf{A}}(\mathrm{H}(\mathbf{C}),\mathrm{H}(\mathbf{C}^{\prime}))\) 其齐次分量的次数为 \(n\) 是 \(\bigoplus_{p+q=n}\operatorname{Tor}_{1}^{\mathbf{A}}(\mathrm{H}_{p}(\mathbf{C}),\mathrm{H}_{q}(\mathbf{C}^{\prime}))\) ，这个连接同态是次数为0的分次同态。此外，存在一个典范同态（X，第62页）

\[
\gamma (\mathrm{B} (\mathrm{C}), \mathrm{C} ^ {\prime}): \mathrm{B} (\mathrm{C}) \otimes_ {\mathrm{A}} \mathrm{H} (\mathrm{C} ^ {\prime}) \rightarrow \mathrm{H} (\mathrm{B} (\mathrm{C}) \otimes_ {\mathrm{A}} \mathrm{C} ^ {\prime}).
\]

在这些记号下：

定理3. --- 假设 A-模 B(C) 和 Z(C) 是平坦的。存在唯一的一个次数为 -1 的分次 k-模同态，

\[
\alpha : \mathrm{H} (\mathrm{C} \otimes_ {\mathrm{A}} \mathrm{C} ^ {\prime}) \rightarrow \operatorname{Tor} _ {1} ^ {\mathrm{A}} (\mathrm{H} (\mathrm{C}), \mathrm{H} (\mathrm{C} ^ {\prime}))
\]

从而使得图表交换

\[
\begin{array}{c} \mathrm{H} (\mathrm{C} \otimes_ {\mathrm{A}} \mathrm{C} ^ {\prime}) \xrightarrow {\alpha} \operatorname{Tor} _ {1} ^ {\mathrm{A}} (\mathrm{H} (\mathrm{C}), \mathrm{H} (\mathrm{C} ^ {\prime})) (- 1) \\ \mathrm{H} (\delta \otimes 1) \Biggl \downarrow \qquad \qquad \qquad \qquad \qquad \qquad \qquad \qquad \qquad \qquad \qquad \qquad \Biggl \downarrow \hat {c} (\Pi , \mathrm{H} (\mathrm{C} ^ {\prime})) \\ \mathrm{H} (\mathrm{B} (\mathrm{C}) \otimes_ {\mathrm{A}} \mathrm{C} ^ {\prime}) (- 1) \xleftarrow {\gamma (\mathrm{B} (\mathrm{C}) . \mathrm{C} ^ {\prime})} (\mathrm{B} (\mathrm{C}) \otimes_ {\mathrm{A}} \mathrm{H} (\mathrm{C} ^ {\prime})) (- 1). \end{array}
\]

分次k模的序列

\[
0 \longrightarrow \mathrm{H} (\mathrm{C}) \otimes_ {\mathrm{A}} \mathrm{H} \left(\mathrm{C} ^ {\prime}\right) \xrightarrow {\gamma \left(\mathrm{C} , \mathrm{C} ^ {\prime}\right)} \mathrm{H} \left(\mathrm{C} ^ {\prime} \otimes_ {\mathrm{A}} \mathrm{C} ^ {\prime}\right) \xrightarrow {\alpha} \operatorname{Tor} _ {1} ^ {\mathrm{A}} \left(\mathrm{H} (\mathrm{C}), \mathrm{H} \left(\mathrm{C} ^ {\prime}\right)\right) (- 1) \longrightarrow 0
\]

是正合的。

因此对每个 \( n \)，我们有一个正合序列

\[
\begin{array}{c} 0 \longrightarrow \bigoplus_ {p + q = n} \mathrm{H} _ {p} (\mathrm{C}) \otimes_ {\mathrm{A}} \mathrm{H} _ {q} (\mathrm{C} ^ {\prime}) \xrightarrow {\gamma_ {n} (\mathrm{C} , \mathrm{C} ^ {\prime})} \mathrm{H} _ {n} (\mathrm{C} \otimes_ {\mathrm{A}} \mathrm{C} ^ {\prime}) \\ \xrightarrow {\alpha_ {n}} \bigoplus_ {p + q = n - 1} \operatorname{Tor} _ {1} ^ {\mathrm{A}} \left(\mathrm{H} _ {p} (\mathrm{C}), \mathrm{H} _ {q} (\mathrm{C} ^ {\prime})\right) \longrightarrow 0. \end{array}\tag{11}
\]

为简单起见，设 \(\mathbf{B} = \mathbf{B}(\mathbf{C})\) , \(\mathbf{Z} = \mathbf{Z}(\mathbf{C})\) , \(\mathbf{H} = \mathbf{H}(\mathbf{C})\) 和 \(\mathbf{H}' = \mathbf{H}(\mathbf{C}')\)：

由于 B 是平坦的，我们从 (I) 推出一个正合序列（X，第 72 页，推论 2）

\[
0 \longrightarrow Z \otimes_ {A} C ^ {\prime} \xrightarrow {j \otimes 1} C \otimes_ {A} C ^ {\prime} \xrightarrow {\delta \otimes 1} (B \otimes_ {A} C ^ {\prime}) (- 1) \longrightarrow 0.\tag{12}
\]

引理 3. --- 与正合序列 (12) 相关联的连接同态 H(B ⊗\_A C') → H(Z ⊗\_A C') 等于 H(i ⊗ 1)。

事实上，设 \(a \in \mathbf{Z}(\mathbf{B} \otimes_{\mathbf{A}} \mathbf{C}')\) ；由于 B 是平坦的， \(a\) 属于 \(\mathbf{B} \otimes_{\mathbf{A}} \mathbf{Z}(\mathbf{C}')\) 的像，因此可以写成 \(\sum_{\lambda} da_{\lambda} \otimes b_{\lambda}\) ，具有 \(a_{\lambda} \in \mathbf{C}\) , \(b_{\lambda} \in \mathbf{C}'\) , \(db_{\lambda} = 0\) 。

\(a\) 的类在所求同态下的像按定义是 \(\mathbf{D}(\sum a_{\lambda}\otimes b_{\lambda})=\sum da_{\lambda}\otimes b_{\lambda}=(i\otimes1)(a)\) 的类，从而引理得证。

因此，与（12）相关联的同调正合列是

\[
\begin{array}{r l} \mathrm{H} (\mathrm{B} \otimes_ {\mathrm{A}} \mathrm{C} ^ {\prime}) & \xrightarrow {\mathrm{H} (i \otimes 1)} \mathrm{H} (\mathrm{Z} \otimes_ {\mathrm{A}} \mathrm{C} ^ {\prime}) \xrightarrow {\mathrm{H} (j \otimes 1)} \mathrm{H} (\mathrm{C} \otimes_ {\mathrm{A}} \mathrm{C} ^ {\prime}) \\ & \xrightarrow {\mathrm{H} (\delta \otimes 1)} \mathrm{H} (\mathrm{B} \otimes_ {\mathrm{A}} \mathrm{C} ^ {\prime}) (- 1) \xrightarrow {\mathrm{H} (i \otimes 1)} \mathrm{H} (\mathrm{Z} \otimes_ {\mathrm{A}} \mathrm{C} ^ {\prime}) (- 1). \end{array}
\]

此外，由于Z是平坦的，我们从(II)推出一个分次k-模的正合序列

\[
0 \longrightarrow \operatorname{Tor} _ {1} ^ {\mathrm{A}} (\mathrm{H}, \mathrm{H} ^ {\prime}) \xrightarrow {\hat {c} (\mathrm{II} , \mathrm{H} ^ {\prime})} \mathrm{B} \otimes_ {\mathrm{A}} \mathrm{H} ^ {\prime} \xrightarrow {i \otimes 1} Z \otimes_ {\mathrm{A}} \mathrm{H} ^ {\prime} \xrightarrow {p \otimes 1} \mathrm{H} \otimes_ {\mathrm{A}} \mathrm{H} ^ {\prime} \longrightarrow 0;
\]

最后，我们有第1号中的典范同态

\[
\gamma_ {\mathbf {B}} = \gamma (\mathbf {B}, \mathbf {C} ^ {\prime}): \mathbf {B} \otimes_ {\mathbf {A}} \mathrm{H} ^ {\prime} \rightarrow \mathrm{H} (\mathbf {B} \otimes_ {\mathbf {A}} \mathbf {C} ^ {\prime})
\]

\[
\gamma_ {Z} = \gamma (Z, C ^ {\prime}): Z \otimes_ {A} H ^ {\prime} \rightarrow H (Z \otimes_ {A} C ^ {\prime})
\]

\[
\gamma_ {\mathrm{C}} = \gamma (\mathrm{C}, \mathrm{C} ^ {\prime}): \mathrm{H} \otimes_ {\mathrm{A}} \mathrm{H} ^ {\prime} \rightarrow \mathrm{H} (\mathrm{C} \otimes_ {\mathrm{A}} \mathrm{C} ^ {\prime}),
\]

由此得到一个分次k-模的图表，其行是正合的

\[
\begin{array}{r l} & \mathbf {B} \otimes \mathbf {H} ^ {\prime} \xrightarrow [ \gamma_ {B} ]{i \otimes 1} \mathbf {Z} \otimes \mathbf {H} ^ {\prime} \xrightarrow [ \gamma_ {Z} ]{p \otimes 1} \mathbf {H} \otimes \mathbf {H} ^ {\prime} \longrightarrow 0 \\ & \mathrm{H} (\mathbf {B} \otimes \mathbf {C} ^ {\prime}) \xrightarrow {\mathrm{H} (i \otimes 1)} \mathrm{H} (\mathbf {Z} \otimes \mathbf {C} ^ {\prime}) \xrightarrow {\mathrm{H} (j \otimes 1)} \mathrm{H} (\mathbf {C} \otimes \mathbf {C} ^ {\prime}) \xrightarrow {\mathrm{H} (\delta \otimes 1)} \mathrm{H} (\mathbf {B} \otimes \mathbf {C} ^ {\prime}) (- 1) \xrightarrow {\mathrm{H} (i \otimes 1)} \mathrm{H} (\mathbf {Z} \otimes \mathbf {C} ^ {\prime}) (- 1) \\ & 0 \longrightarrow \operatorname{Tor} _ {1} ^ {\mathrm{A}} (\mathrm{H}, \mathrm{H} ^ {\prime}) (- 1) \xrightarrow {\delta (\mathrm{II.H} ^ {\prime})} (\mathbf {B} \otimes \mathbf {H} ^ {\prime}) (- 1) \xrightarrow [ i \otimes 1 ]{\gamma_ {B}} (\mathbf {Z} \otimes \mathbf {H} ^ {\prime}) (- 1), \end{array}
\]

这由同态 \(\gamma\) 的定义是交换的。但是，由于复形 B 和 Z 是分裂且平坦的， \(\gamma_{\mathbf{B}}\) 和 \(\gamma_{\mathbf{Z}}\) 是双射的（X，第 65 页，推论 1）。由此一方面可推出， \(\gamma_{\mathbf{C}}\) 是单射，其像等于 Ker H( \(\delta \otimes 1\) )，另一方面， \(\gamma_{\mathbf{B}} \circ \partial(\mathrm{II}, \mathrm{H}')\) 是单射，其像等于 Im H( \(\delta \otimes 1\) )。定理由此直接得出。

推论1. --- 若B(C)和Z(C)是平坦的，则对每个左A模N和每个整数n，存在一个正合序列

\[
0 \longrightarrow \mathrm{H} _ {n} (\mathrm{C}) \otimes_ {\mathrm{A}} \mathrm{N} \xrightarrow {\gamma_ {n}} \mathrm{H} _ {n} (\mathrm{C} \otimes_ {\mathrm{A}} \mathrm{N}) \xrightarrow {\alpha_ {n}} \operatorname{Tor} _ {1} ^ {\mathrm{A}} \left(\mathrm{H} _ {n - 1} (\mathrm{C}), \mathrm{N}\right) \longrightarrow 0.\tag{13}
\]

推论2.——设 B(C) 和 B(C') 是投射的，且 Z(C) 是平坦的。则 k-模序列 (11) 和 (13) 是正合且分裂的。

这由定理及以下引理得出：

引理4. --- 若 B(C) 和 B(C') 是投射的，则典范同态

\[
\gamma (\mathbf {C}, \mathbf {C} ^ {\prime}): \mathrm{H} (\mathbf {C}) \otimes_ {\mathbf {A}} \mathrm{H} (\mathbf {C} ^ {\prime}) \rightarrow \mathrm{H} (\mathbf {C} \otimes_ {\mathbf {A}} \mathbf {C} ^ {\prime})
\]

有一个 k-线性收缩。

事实上，根据 X，第 65 页，注记 b)，存在拟同构 \(\varphi: C \to H(C)\) 和 \(\varphi': C' \to H(C')\) 使得 \(H(\varphi) = 1_{H(C)}\) 和 \(H(\varphi') = 1_{H(C')}\) 。在交换图中

\[
\begin{array}{c} \mathrm{H(C)} \otimes_ {\mathrm{A}} \mathrm {H(C^ {\prime})} \xrightarrow {\gamma (\mathrm{C,C} ^ {\prime})} \mathrm {H(C} \otimes_ {\mathrm{A}} \mathrm {C^ {\prime})} \\ \mathrm {H}(\varphi) \otimes \mathrm {H}(\varphi^ {\prime}) \Big \downarrow \\ \mathrm {H(C)} \otimes_ {\mathrm{A}} \mathrm {H(C^ {\prime})} \xrightarrow {\gamma (\mathrm{H(C),H(C^ {\prime})})} \mathrm {H(C)} \otimes_ {\mathrm{A}} \mathrm {H(C^ {\prime})} \end{array}
\]

\(\mathrm{H}(\varphi) \otimes \mathrm{H}(\varphi') \text{ 和 } \gamma(\mathrm{H}(\mathrm{C}), \mathrm{H}(\mathrm{C}') ) \text{ 是恒等映射，由此得出断言。}\)

推论3（“泛系数公式”）。—— 设环 A 是主理想环。若复形 C 与 C' 是自由的，则 A-模序列 (11) 是正合且分裂的；若复形 C 是自由的，则对每个 A-模 N，A-模序列 (13) 是正合且分裂的。

事实上，B(C)、Z(C) 和 B(C') 是自由模 C、C、C' 的子模，因此是自由的（VII，§3，定理1的推论2），于是应用推论2。

推论4（“屈内特公式”）。——设 C 上有界，且 C 与 H(C) 是平坦的；则典范同态

\[
\gamma (\mathrm{C}, \mathrm{C} ^ {\prime}): \mathrm{H} (\mathrm{C}) \otimes_ {\mathrm{A}} \mathrm{H} (\mathrm{C} ^ {\prime}) \rightarrow \mathrm{H} (\mathrm{C} \otimes_ {\mathrm{A}} \mathrm{C} ^ {\prime})
\]

是双射。

由该定理，只需证明 B(C) 和 Z(C) 是平坦的。现在我们有以下正合序列

\[
0 \to \mathrm{B} _ {n} (\mathrm{C}) \to \mathrm{Z} _ {n} (\mathrm{C}) \to \mathrm{H} _ {n} (\mathrm{C}) \to 0
\]

\[
0 \to Z _ {n} (\mathbf {C}) \to \mathbf {C} _ {n} \to \mathbf {B} _ {n - 1} (\mathbf {C}) \to 0,
\]

由此，根据 X 第 75 页推论 1，可得如下蕴含关系 \((\mathbf{B}_{n-1}(\mathbf{C})\) 是平坦的) \(\Rightarrow (\mathbf{Z}_n(\mathbf{C})\) 是平坦的) \(\Rightarrow (\mathbf{B}_n(\mathbf{C})\) 是平坦的)；我们注意到 \(\mathbf{B}_n(\mathbf{C}) = 0\) 对足够小的 \(n\) 成立，由此得出结论。

推论5. --- 设 \(u: \mathbf{C} \to \mathbf{C}'\) 为右 A-模复形的同态，平坦且右有界。对于左 A-模的每个复形 E，态射 \(u \otimes 1_{\mathbf{E}}: \mathbf{C} \otimes_{\mathbf{A}} \mathbf{E} \to \mathbf{C}' \otimes_{\mathbf{A}} \mathbf{E}\) 是拟同构。

确实，Con \((u)\) 是一个平坦复形，右有界且同调为零；因此由推论4，我们有 H(Con \((u) \otimes_{\mathbf{A}} \mathbf{E}) = 0\) ，故 H(Con \((u \otimes 1_{\mathbf{E}})) = 0\) （X，第 67 页，引理2），从而 \(u \otimes 1_{\mathbf{E}}\) 是拟同构。

